% !PRINTED 649
\weberpage{649}
%<<CONT
then $\mathfrak{G}\,|\,\mathfrak{A}$ (\S.\,4 of this volume) is the group of the field $\Omega_1$, and,
like $\mathfrak{G}$ itself, this too is an \gesperrt{Abelian} group. A divisor
of $\Omega$ which is of the same degree as $\Omega$ is identical with $\Omega$
(cf. Bd. I, \S.\,144, 149, 156).

If now there are two non-primitive numbers $x_1, x_2$ in $\Omega$ of
such a kind that we can set
\[
\Omega = R(x_1, x_2)
\]
then the fields $\Omega_1 = R(x_1)$, $\Omega_2 = R(x_2)$ are
proper divisors of $\Omega$, and $\Omega$ is said to be \gesperrt{composed} of the two fields
$\Omega_1, \Omega_2$.

If the field $\Omega$ admits no such representation, it is called
\gesperrt{simple}.

If, then, it is assumed as proved that $\Omega_1, \Omega_2$ are
cyclotomic fields, i.\,e. that all their numbers can be represented
rationally by roots of unity, the same follows for $\Omega$.

If one of the fields $\Omega_1, \Omega_2$ is decomposable, we can
repeat the decomposition, and since the degrees are decreasing
positive integers, we must finally arrive at simple
fields.

\gesperrt{Theorem I thus needs to be proved only for simple Abelian
fields.}

It is to be remembered, however, that the decomposability of a
field $\Omega$ by no means follows merely from the existence of a
divisor, and that for a decomposable field the decomposition into
simple fields can take place in quite different ways, so
that for fields the laws of divisibility do not hold
as in the region of numbers.

A criterion for the simplicity of a field we can
derive from its group.

If the group $\mathfrak{G}$ has two proper divisors $\mathfrak{A}$ and $\mathfrak{B}$ of such a
kind that every element can be composed once and only once from an
element of $\mathfrak{A}$ and an element of $\mathfrak{B}$, we call the group
$\mathfrak{G}$ decomposable into the two
components $\mathfrak{A}, \mathfrak{B}$; $\mathfrak{G}$ is thus decomposable if, in the product
\weq{1}{\mathfrak{G} = \mathfrak{A}\mathfrak{B}}
formed according to the composition of the parts (\S.\,4), every
element of $\mathfrak{G}$ appears once and only once. The degree
of $\mathfrak{G}$ is then equal to the product of the degrees of $\mathfrak{A}$ and $\mathfrak{B}$.

If there are no such divisors, the group $\mathfrak{G}$ is called
indecomposable.
